\section{Conclusion}\label{sec:conclusion}

This thesis asked whether a general-purpose zero-knowledge virtual machine (zkVM) is a viable substrate for post-quantum anonymous credentials on consumer hardware. The goal was to make BDEC credential verification, instantiated with the PLUM signature, generate proofs in a practically acceptable time inside a zkVM on a personal computer (Apple M5~Pro, $24$~GB); the means was a custom precompile suite that absorbs the field-mismatch tax of PLUM's large-prime-field algebraic hashing. We realised the construction on RISC~Zero and SP1, measured it against an Aurora static-circuit baseline, and analysed what the substitution preserves and what it costs.

\subsection{Main Findings}\label{subsec:conclusion-findings}

The central result is a frontier, and neither substrate is a clear winner. On the target hardware the BDEC proof one can afford to generate is not zero-knowledge, and the zero-knowledge wrap the toolchain exposes is not post-quantum (Section~\ref{ssec:zk-gap}). The first prong is a measured resource bound: the PLONK wrap of the standalone PLUM verification receipt is killed under memory pressure at ${\approx}20$~GB of the $24$~GB budget. The second is a fact about the deployed toolchain: the only succinct wraps it offers are pairing-based, so additional memory would not restore post-quantum anonymity. This obstruction on both sides holds independently of any timing number, and the rest of the thesis rests on it.

Around that finding, the measurements establish the surrounding cost picture. The custom Griffin-$\mathbb{F}_p^{192}$ precompile turns a non-terminating emulation into a finite measurement: without it, standalone PLUM verification exhausts memory at ${\approx}1$m$45$s on the target hardware; with it, the verification proves in $32.5$ minutes on SP1 at $\lambda=80$. Two of the measurements falsify natural predictions. First, the bespoke algebraic-hash precompile is \emph{slower} than a software SHA-3 control at matched security ($32.5$ versus $13.3$ minutes in prove wall-clock, about $2.45\times$): arithmetising a $199$-bit-field algebraic hash through a custom chip costs more than a hash the substrate already supports. Second, against the transparent, non-preprocessing Aurora baseline the zkVM's update-churn advantage never shows up in wall-clock. The static-side recompile is sub-second ($R_{\mathsf{static}}\approx 0.3$~s, Section~\ref{ssec:churn-eval}), so what separates the two regimes is the regeneration-and-re-audit \emph{footprint} (Table~\ref{tab:churn}), a cost in engineering and assurance rather than in proving time. A wall-clock flexibility advantage would appear only against a \emph{preprocessing} baseline, which is not what BDEC uses.

The BDEC relations themselves run correctly inside the zkVM: credential generation executes both of its signature verifications to acceptance in $9.4\times10^{9}$ guest cycles, and the presentation relation in $1.41\times10^{10}$ (one credential shown) and $1.88\times10^{10}$ (two) cycles, consistent with the additive cost model. No zero-knowledge \emph{proof} of either relation was produced on consumer hardware: the credential-generation succinct prove terminated in an internal-verifier rejection (RISC~Zero~$3.0.3$; an anomaly under diagnosis, categorically distinct from a resource limit and on which no claim of this thesis rests, Section~\ref{ssec:cregen-anomaly}), and the presentation wrap is bounded by the same memory frontier.

\subsection{Answer to the Research Question}\label{subsec:conclusion-answer}

The answer is \emph{not yet}, for a reason more specific than ``too slow.'' Standalone PLUM verification becomes feasible on consumer hardware once the precompile absorbs the field-mismatch tax. Deployable post-quantum \emph{anonymity} is not reached, because it requires a zero-knowledge wrap that is at once feasible to generate within $24$~GB and post-quantum-sound, and the current toolchains provide neither. The substrate choice is therefore governed by the deployment-lifecycle decision rule of Section~\ref{sec:discussion}: where the verification predicate is fixed and proofs are frequent, the specialised static circuit is preferable; where the predicate, signature, or security parameters evolve over the deployment's lifetime, the zkVM's fixed universal relation confines each change to a guest-program edit, provided the change leaves the precompiled primitive suite intact, since a hash-family or field-width change regenerates the Griffin AIR. Against a non-preprocessing baseline its advantage there is a smaller re-audit footprint, not a faster runtime.

\subsection{Limitations}\label{subsec:conclusion-limitations}

The results are bounded in ways we state explicitly. Prove-mode wall-clock is measured at $\lambda=80$, the largest level at which it is observable within the memory budget and not a deployment security level; cycle attribution is reported at $\lambda=128$. The static baseline is Loquat over $\mathbb{F}_{p_{127}}$ with zero-knowledge disabled, not same-scheme PLUM, so the per-proof comparison is indicative and the lifecycle crossover is bounded in form but not located in wall-clock. The prove-mode figures are single runs. The security guarantees are conditional: the constraint-level soundness of the Griffin AIR (premise~(T2)), the existence of a PLUM signature-transcript simulator (premise~(T4)), the quantum-random-oracle lifts, and the canonical prime instantiation are carried as open premises (Section~\ref{sec:security}). Finally, the RISC~Zero prototype does not yet bind the receipt journal to the public statement, so the artifact is at present a functional benchmark rather than a statement-bound credential verifier.

\subsection{Future Work}\label{subsec:conclusion-future}

The most informative next measurement is a \emph{same-scheme} per-proof comparison, either running Loquat verification inside the zkVM to match the static Loquat/Aurora point or building an $\mathbb{F}_p$ Aurora harness to run PLUM statically; this would locate the crossover that the decision rule currently bounds only in form. The open premises are the substantive cryptographic work that remains: premise~(T2) by metamorphic and fault-injection testing of the AIR, or constraint-level verification of the chip against the reference permutation, by the method of Arguzz~\cite{hochrainer2025arguzz}; premise~(T4) by constructing a PLUM signature-transcript simulator; and the post-quantum-anonymity goal by a post-quantum-sound zero-knowledge wrap (a STARK- or lattice-based recursive argument) in place of the pairing-based wraps the toolchains currently expose. The immediate engineering step is to commit the public statement to the receipt journal, after which the credential-generation prove anomaly should be re-run with segment-level diagnostics. Confirming the canonical prime with PLUM's authors would discharge the prime-instantiation premise.

\subsection{Closing}\label{subsec:conclusion-closing}

Future-proof post-quantum identity systems require both cryptographic soundness and operational evolvability, and this thesis measured the price of pursuing both on the same substrate. What it contributes is a mapped frontier. The precompile makes post-quantum signature verification finite to prove on a personal computer. The deployment-lifecycle separation between the static and zkVM regimes is real, though it surfaces as a re-audit footprint and stays out of the wall-clock numbers. The obstruction to deployable post-quantum anonymity has two prongs that survive together: the proof one can afford is feasible but not anonymous, and the anonymous wrap is not post-quantum, and no amount of additional memory or favourable timing removes either. That obstruction, named precisely and surrounded by its cost picture, is what we hand to the design of the next such system.
